\documentclass{beamer}
\mode<presentation>
\usepackage{amsmath}
\usepackage{amssymb}
%\usepackage{advdate}
\usepackage{graphicx}
\graphicspath{{../figs/}}
\usepackage{adjustbox}
\usepackage{subcaption}
\usepackage{enumitem}
\usepackage{multicol}
\usepackage{mathtools}
\usepackage{listings}
\usepackage{url}
\def\UrlBreaks{\do\/\do-}
\usetheme{Boadilla}
\usecolortheme{lily}
\setbeamertemplate{footline}
{
  \leavevmode%
  \hbox{%
  \begin{beamercolorbox}[wd=\paperwidth,ht=2.25ex,dp=1ex,right]{author in head/foot}%
    \insertframenumber{} / \inserttotalframenumber\hspace*{2ex} 
  \end{beamercolorbox}}%
  \vskip0pt%
}
\setbeamertemplate{navigation symbols}{}
\let\solution\relax
\usepackage{gvv}
\lstset{
%language=C,
frame=single, 
breaklines=true,
columns=fullflexible
}

\numberwithin{equation}{section}
\title{12.287}
\author{AI25BTECH11027 - NAGA BHUVANA}
\begin{document}
{\let\newpage\relax\maketitle}
\renewcommand{\thefigure}{\theenumi}
\renewcommand{\thetable}{\theenumi}

\noindent
		\textbf{Question:}\\
An Unitary matrix   $\myvec{ae^{i\alpha} & b\\ ce^{i\beta} & d}$ is given where $a,b,c,d,\alpha,\beta$ are real. The inverse of the matrix is\\
    a) $\myvec{ae^{i\alpha} & -ce^{i\beta}\\ b & d}$ 
    b) $\myvec{ae^{i\alpha} & ce^{i\beta}\\ b & d}$ 
    c) $\myvec{ae^{-i\alpha} & b\\ ce^{-i\beta} & d}$ 
    d) $\myvec{ae^{-i\alpha} & ce^{-i\beta}\\ b & d}$ 
\textbf{Solution:}\\
\textbf{Unitary Matrix:} Complex square matrix whose inverse is equal to its conjugate transpose\\
Let \begin{align}
    M=\myvec{ae^{i\alpha} & b\\ ce^{i\beta} & d}
\end{align}
$\bar{M}$ be Complex Conjugate Matrix of $M$\\
\begin{align}
    \bar{M}=\myvec{ae^{-i\alpha} & b\\ ce^{-i\beta} & d}
\end{align}
\begin{align}
    M^{-1}=\bar{M}^{T}
\end{align}
\begin{align}
    M^{-1}=\myvec{ae^{-i\alpha} & b\\ ce^{-i\beta} & d}^T
\end{align}
\begin{align}
    M^{-1}=\myvec{ae^{-i\alpha} & ce^{-i\beta}\\ b & d}
\end{align}
$\therefore$ Inverse of the Matrix is $\myvec{ae^{-i\alpha} & ce^{-i\beta}\\ b & d}$
\end{document}
